The conducted \abbr{SWOT} analysis indicates that the proposed satellite-based Early Warning Alarm System occupies a strong strategic position due to its unique operational capabilities and its alignment with the growing demand for civil protection and national resilience in the Baltic region.

\subsection{Strengths}

The most significant strengths identified are the high reliability of the satellite constellation and the absence of a dedicated regional early-warning capability offering comparable functionality. These factors achieved the highest weighted scores, demonstrating that the proposed system addresses a critical capability gap while providing a reliable and innovative solution for missile early warning. Early threat detection and high system availability further enhance the mission's operational value, while the capability to store historical event data supports post-event analysis, continuous improvement of detection algorithms, reduction of false alarms, and refinement of future warning products.

\subsection{Opportunities}

The identified opportunities strongly complement these internal strengths. The most important external factors include the increasing demand for effective civil protection and missile early-warning capabilities in the Baltic region, together with the availability of European Union and \abbr{NATO} programmes supporting resilience, crisis management, and dual-use technologies.

The interaction analysis demonstrates that the project's strengths effectively enable the exploitation of these opportunities. In particular, the combination of high system reliability, continuous monitoring capability, and early-warning functionality creates favourable conditions for institutional cooperation, long-term programme support, and integration into future regional civil protection and resilience frameworks. Furthermore, the potential expansion of the satellite constellation provides additional long-term development opportunities, although its implementation will depend on technical feasibility and the availability of financial resources.

\subsection{Weaknesses}

Despite these advantages, several internal weaknesses may reduce the project's ability to fully capitalize on emerging opportunities. The most significant weakness is the high cost of satellite deployment and long-term operations, which directly affects the feasibility of future constellation expansion and the development of the mission's ground segment.

Additional challenges include gradual satellite degradation throughout the mission lifetime, limited maintenance possibilities in orbit, and the complexity of integrating observations from multiple satellites into a coherent monitoring system. False-positive detections caused by atmospheric phenomena may also reduce operational efficiency if filtering and classification methods are not continuously improved. Overall, these weaknesses primarily affect the scalability and long-term sustainability of the system rather than its core operational objectives.

\subsection{Threats}

The analysis also identifies several external threats that could negatively influence the long-term success of the mission. The most significant risks are cyberattacks and anti-satellite activities, including jamming and spoofing, which received the highest weighted importance among the identified threats.

Since the mission relies on continuous satellite operation, secure communications, and reliable data transmission between the space and ground segments, maintaining robust cybersecurity and resilience against electronic interference is essential. An additional critical vulnerability is the dependence on Inter-Satellite Links connecting the \abbr{LEO} tracking satellites with the \abbr{GEO} communication and early-warning satellites. Disruption, jamming, or spoofing of these links could significantly degrade the system's ability to transmit detection data in real time, resulting in delayed or interrupted dissemination of warning information.

Additional threats include political or legal restrictions affecting international information sharing, the increasing risk of collisions with space debris, and dependence on foreign launch providers, which may introduce schedule uncertainties and supply chain vulnerabilities.

The interaction analysis between weaknesses and threats highlights that operational limitations---including limited maintenance capabilities, satellite degradation, the complexity of multi-satellite data integration, and reliance on uninterrupted inter-satellite communication links---significantly increase the system's exposure to external risks. These factors may reduce mission resilience in the event of cyber incidents, hostile electronic interference, \abbr{ISL} disruption, or failures caused by the orbital environment. Consequently, mitigating these weaknesses through appropriate system design, communication redundancy, robust cybersecurity measures, reliable operational procedures, and careful mission planning should remain a priority throughout the system development process.

\subsection{Strategy}

The overall strategic assessment clearly indicates that the project should adopt an \textbf{Attack (\abbr{SO}) Strategy}. This strategy achieved the highest overall score, with the greatest number of interactions and the highest weighted value among all strategic options. The results indicate that the project's strengths significantly outweigh its weaknesses and can be effectively leveraged to exploit favourable external opportunities.

The system's unique regional capability, combined with the growing demand for missile early-warning systems supporting civil protection and the availability of institutional funding programmes, creates favourable conditions for successful implementation, future expansion, and long-term operational deployment.

Although the \textbf{Retreat (\abbr{WT}) Strategy} achieved the second-highest score, reflecting the presence of significant technical and operational risks, these challenges do not outweigh the project's strategic advantages. Instead, they highlight the importance of proactive risk mitigation, particularly in the areas of cybersecurity, satellite reliability, launch independence, communication resilience, and long-term system maintainability.

\subsection{Conclusions}

The \abbr{SWOT} analysis demonstrates that the proposed hybrid \abbr{GEO}-\abbr{LEO} satellite constellation represents a promising and strategically justified investment in regional civil protection and missile early-warning capabilities. Its principal advantages-including early threat detection, high operational reliability, continuous availability, and the ability to archive and analyse historical event data-provide a strong foundation for addressing the growing resilience requirements of the Baltic region.

Nevertheless, the long-term success and operational effectiveness of the \abbr{EWAS} mission will depend on the effective management of technical, operational, and security-related risks. Particular attention should be devoted to strengthening cybersecurity, improving communication resilience, ensuring system redundancy, and maintaining the long-term sustainability of the satellite constellation. Addressing these challenges during the design, implementation, and operational phases will significantly increase the mission's reliability and overall strategic value.